% ===================================================================
%  اللوحة رقم ١٢ — المحطّة. --- سكك الحديد. --- السفن.
%
%  Traduction, faite en 2026, en arabe levantin d'aujourd'hui, sur
%  l'ido de texto/io. Regles de la colonne : voir 10-lawha-01.tex.
%  Une seule numerotation. Le « Noto. » d'ouverture reste dans le
%  bloc apar, sans cle a lui ; la note de bas de page garde la
%  sienne, t12-noto-1, posee apres le bloc 14 — que l'ido coupe en
%  deux par un « suite » et que la traduction ecrit d'un seul tenant.
%
%  LE TELEGRAMME (5) EST « برقية », UN ECLAIR. Le mot a ete forge en
%  arabe au XIXe siecle sur برق, la foudre : une برقية, c'est un
%  coup de foudre expedie. La colonne gujaratie avait « તાર », le
%  fil. DEUX LANGUES ONT NOMME LA MEME INVENTION AU MEME MOMENT,
%  L'UNE PAR LE FIL QU'ON TEND, L'AUTRE PAR L'ECLAIR QUI PASSE
%  DEDANS — et le bureau (4) est resté « مكتب التلغراف », transcrit,
%  parce que le batiment est venu tout fait.
%
%  L'HORLOGE DE LA FRONTIERE (89) AVANCE DE CINQUANTE-CINQ MINUTES,
%  ET CETTE COLONNE-CI N'A PAS BESOIN DE FRONTIERE POUR CONNAITRE LA
%  SCENE. La colonne gujaratie avait rappelé l'heure de Bombay,
%  distincte de l'heure indienne jusqu'en 1955. Ici c'est plus simple
%  et plus ancien : jusque dans les annees 1920, le Levant comptait
%  les heures A LA TURQUE, a partir du coucher du soleil, en meme
%  temps que l'heure europeenne dite « a la franque », et les
%  horloges publiques de Damas et de Beyrouth portaient les deux.
%  DEUX HEURES SUR LE MEME CADRAN, SANS QU'ON AIT A CHANGER DE PAYS.
%
%  LE CHEMIN DE FER DE CE PAYS-LA A UNE DATE ET DEUX ORIGINES, ET LE
%  VOCABULAIRE LES PORTE. La ligne Beyrouth-Damas est francaise et
%  date de 1895 ; la ligne du Hedjaz, de Damas vers Medine, est
%  ottomane et date de 1900, batie par souscription. Le lexique tient
%  des deux : « محطّة » la gare, « قطار » le train, « قاطرة » la
%  locomotive (54), « تذكرة » le billet (9) sont arabes ; « رصيف »
%  le quai est arabe aussi ; mais « البوفيه » (34) et
%  « الجرسون » (35) sont francais, et « الشبّاك » du guichet est le
%  mot ordinaire de la fenetre. CE QU'ON ACHETE ET CE QU'ON PREND EST
%  NOMME ICI ; CE QU'ON Y CONSOMME EST NOMME AILLEURS.
%
%  Enfin les bagages (17) sont « عفش » et le porteur (16)
%  « العتّال », deja rencontre au tableau 5 sur le chantier : le meme
%  homme, le meme mot, une charge sur le dos. La colonne gujaratie
%  avait « કુલી » ici, et c'etait le seul mot du releve parti d'Inde
%  vers l'anglais. Celle-ci n'a rien prete a ce renvoi-la — elle
%  avait deja prete l'arsenal et l'amiral deux tableaux plus tot.
% ===================================================================

%%K t12-apar-1 apar
\VUsaut{4.0mm}
\VUtitre{15.0pt}{60}{اللوحة رقم ١٢}
\VUsaut{2.4mm}
\VUpk{11.4pt}{40}{المحطّة. --- سكك الحديد. --- السفن.}
\VUsaut{3.4mm}
\textsc{ملاحظة.} --- جداول المواعيد المستعملة بكل بلد مختلفة، لهيك
منستعمل كمثال ساعات \VUgras{الجدول الفرنسي}.

%%K t12-01-1 p
1. --- هلّق رجعت من سفرة بألمانيا. وقبل ما سافر، اشتريت \VUgras{دليل
المواعيد} \textsuperscript{(15)} لأعرف ساعات سفر القطارات. وبعد ما
تأكّدت إنّه أنسب قطار رح يسافر من باريس الساعة تسعة بالليل ورح يوصل
عسترازبورغ الساعة وحدة وتلتعش دقيقة، حضّرت سفرتي، وتاني يوم خلّيتهن
يوصّلوني عالمحطّة.

%%K t12-02-1 p
2. --- ولمّا فتّ عقاعة السفر الكبيرة، وراء \VUgras{خوري}
\textsuperscript{(2)} ضيعة عتيق كان شايل بإيده \VUgras{شنتا سفره}
\textsuperscript{(3)}، كان في \VUgras{مسافرين} \textsuperscript{(1)} كتار
واصلين من قبل.

%%K t12-02-2 p
ولأني كان بدّي ابعت \VUgras{برقية (تلغراف)} \textsuperscript{(5)}،
خلّيت \VUgras{مراقب} \textsuperscript{(6)} يدلّني على \VUgras{مكتب
التلغراف} \textsuperscript{(4)}.

%%K t12-03-1 p
3. --- وقبل ما فوت على \VUgras{قاعة الانتظار} \textsuperscript{(7)}،
رحت جبت \VUgras{تذكرة} \textsuperscript{(9)} ذهاب وإياب.

%%K t12-04-1 p
4. --- وما استنّيت كتير عند \VUgras{الشبّاك} \textsuperscript{(8)}، لأنه
ما كان في ناس كتار هونيك. و\VUgras{عسكري} \textsuperscript{(10)}، كان
مسافر بإجازة، تكرّم وتنازل لي عن دوره، وهيك صار عندي وقت كافي لسأل كم
سؤال \VUgras{ناظر المحطّة} \textsuperscript{(13)}، اللي كان عم يحكي مع
\VUgras{مفوّض} \textsuperscript{(12)} المراقبة ومع \VUgras{الدركي}
\textsuperscript{(11)} المناوب.

%%K t12-tit-1 sub
\VUpk{10.2pt}{40}{تسجيل العفش.}
\VUsaut{3.0mm}

%%K t12-05-1 p
5. --- وبعد ما اشتريت جريدة من \VUgras{كشك الكتب} \textsuperscript{(14)}،
خلّيتهن يوزنوا \VUgras{عفشي} \textsuperscript{(17)} اللي كان هلّق جابه
\VUgras{عتّال} \textsuperscript{(16)} على \VUgras{عربايّة العفش}
\textsuperscript{(19)} ؛ و\VUgras{الموظّف} \textsuperscript{(22)} المكلّف
بـ\VUgras{الميزان} \textsuperscript{(21)} خبّرني إنّه صندوقي وزنه خمسة
وتلاتين كيلو ؛ يعني \VUgras{وزن زيادة} خمس كيلوات، لازم ادفع عنه زيادة
عند \VUgras{شبّاك العفش} \textsuperscript{(23)}.

%%K t12-06-1 p
6. --- ولأني كنت بكّير كتير، صار عندي وقت اتفرّج على حركة المسافرين
بالمحطّة. بعضهن عم يتركوا أو ياخدوا عفشهن الخفيف من \VUgras{مستودع
الأمانات} \textsuperscript{(25)} ؛ وغيرهن عم يتطلّعوا بـ\VUgras{جداول
المواعيد} \textsuperscript{(26)}. والمسافرين الواصلين عم يستعجلوا
عالـ\VUgras{مخرج} \textsuperscript{(32)}، و\VUgras{شنتاهن}
\textsuperscript{(24)} بإيدهن. وبعضهن عم يستنّوا عند \VUgras{شبّاك تسليم
العفش} \textsuperscript{(27)} وعم يقدّموا \VUgras{وصلهن}
\textsuperscript{(29)} ليستلموا صناديقهن، و\VUgras{الصناديق}
\textsuperscript{(28)} أو \VUgras{القفف} \textsuperscript{(30)}.

%%K t12-07-1 p
7. --- و\VUgras{موظّفي المكس} \textsuperscript{(33)} عم يفتّشوا الطرود
بعناية ليتأكّدوا إذا في شي لازم ينصرّح عنه.

%%K t12-08-1 p
8. --- وبعض المسافرين راحوا عالـ\VUgras{بوفيه} \textsuperscript{(34)}،
وين \VUgras{الجرسون} \textsuperscript{(35)} عم يهتمّ بخدمتهن. ولازم
يستعجلوا، لأنه \VUgras{القطار} \textsuperscript{(36)}، اللي هوّي سريع، ما
رح يضلّ كتير بالمحطّة.

%%K t12-09-1 p
9. --- وبعدين رحت عرصيف السفر ودوّرت على \VUgras{عربة}
\textsuperscript{(37)} مباشرة لأني ما بدّي اضطرّ غيّر العربة بالطريق.
وطلعت على عربة فيها ممرّ، فيها \VUgras{مقصورة} \textsuperscript{(38)}
للمدخّنين ومقصورة محجوزة لـ« السيّدات لحالهن ».

%%K t12-tit-2 sub
\VUpk{10.2pt}{40}{السفر بالليل.}
\VUsaut{3.0mm}

%%K t12-10-1 p
10. --- وبعد ما طلعت \VUgras{الدرجة} \textsuperscript{(40)}، وسكّرت
\VUgras{الباب} \textsuperscript{(39)}، ولفّيت \VUgras{الستارة}
\textsuperscript{(41)}، وحطّيت عفشي الخفيف على \VUgras{شبكة العفش}
\textsuperscript{(43)}، قعدت عالـ\VUgras{مقعد} \textsuperscript{(42)}. ولمّا
شفت \VUgras{عامل اللمبات} \textsuperscript{(44)} عم يولّع اللمبات، فكّرت
إنّه لازم قضّي ليلة بالعربة، وناديت الموظّف المكلّف بتأجير
\VUgras{المساند} \textsuperscript{(45)} لأطلب منه واحد.

%%K t12-11-1 p
11. --- وإجا المراقب يتأكّد من التذاكر. وسألته ليش لسّا ما سافرنا،
لأنه صار الوقت المظبوط. وجاوبني إنّه \VUgras{رئيس القطار}
\textsuperscript{(46)} مشغول عم يحطّ بسكليتات بـ\VUgras{عربة الأمتعة}
\textsuperscript{(47)}. وكمان، لسّا ما بدّلوا كل \VUgras{مدافئ الأرجل}
\textsuperscript{(48)}.

%%K t12-12-1 p
12. --- بس كنّا رح نسافر عن قريب، لأنه \VUgras{الوقّاد}
\textsuperscript{(50)}، وهوّي على \VUgras{عربة الفحم} \textsuperscript{(49)}،
أخد فحم ليشعل نار \VUgras{القاطرة} \textsuperscript{(54)}،
و\VUgras{السائق} \textsuperscript{(51)} كان عم يستنّى بس إشارة \VUgras{نايب
ناظر المحطّة} \textsuperscript{(52)} اللي كان عالـ\VUgras{رصيف}
\textsuperscript{(53)}.

%%K t12-13-1 p
13. --- و\VUgras{مرجل} \textsuperscript{(56)} القاطرة كان عم يهدر بخفوت ؛
و\VUgras{المدخنة} \textsuperscript{(55)} كانت عم تقذف سيول دخان مخلوطة
بـ\VUgras{بخار} \textsuperscript{(60)}. و\VUgras{صفّارة}
\textsuperscript{(59)} حادّة صرخت و\VUgras{المكابس} \textsuperscript{(58)}
بلّشت تتحرّك وعم تدفع \VUgras{دواليب} \textsuperscript{(57)} المكنة
التقيلة.

%%K t12-tit-3 sub
\VUsaut{3.0mm}
\VUpk{10.2pt}{40}{سافرنا !}
\VUsaut{3.0mm}

%%K t12-14-1 p
14. --- وسرعة القطار اللي عم يركض عالـ\VUgras{قضبان}
\textsuperscript{(64)} زادت ؛ و\VUgras{الأرصفة} \textsuperscript{(65)}
اختفت ؛ ومرقنا جنب \VUgras{المراحيض} \textsuperscript{(66)}، وكمان جنب
صفّ عربات موقوفة عالـ\VUgras{خطّ} \textsuperscript{(63)} التاني :
\VUgras{عربات نوم} \textsuperscript{(67)}، و\VUgras{عربة طعام}
\textsuperscript{(68)}، و\VUgras{عربة بريد} \textsuperscript{(69)}.

%%K t12-noto-1 noto
\VUnotes{143.2mm}{%
\textit{(*) ملاحظة.} --- \textit{طبعاً وصف هالسفرة حسب الحدود اللي كانت
قبل الحرب.}%
}

%%K t12-15-1 p
15. --- وبعدين مرقت قدّام عيوننا \VUgras{محطّة البضايع}
\textsuperscript{(71)}. وتحت العنابر، الموظّفين عم يحمّلوا
\VUgras{البضايع} \textsuperscript{(72)} اللي رح تنبعت بقطار بطيء.
و\VUgras{فرق} \textsuperscript{(76)} جنب \VUgras{خزّان الميّ}
\textsuperscript{(73)} عم يناوروا \VUgras{عربات المواشي}
\textsuperscript{(75)} و\VUgras{العربات المسطّحة} \textsuperscript{(79)}
ليعملوا \VUgras{قطار بضايع} \textsuperscript{(80)}.

%%K t12-16-1 p
16. --- وقطعنا آخر \VUgras{مفرق} \textsuperscript{(78)}، اللي جنبه كان
واقف عامل التحويلة ؛ ومرقنا جنب \VUgras{قرص الإشارة}
\textsuperscript{(86)} والمحطّة اختفت بعيد.

%%K t12-17-1 p
17. --- وعن قريب قطعنا \VUgras{جسر حديد} \textsuperscript{(74)} عم يمرق
فوق \VUgras{النهر} \textsuperscript{(81)}. وبعد ما قطعنا هالجسر، مشينا
محاذيين النهر، اللي عليه \VUgras{قاطرات} \textsuperscript{(82)} قوية عم
تجرّ \VUgras{صنادل} \textsuperscript{(83)} تقيلة. وكمان شفنا \VUgras{باخرة
ركّاب} \textsuperscript{(84)} وهيّي عم ترسي على \VUgras{الرصيف العايم}
\textsuperscript{(85)}.

%%K t12-18-1 p
18. --- وبعد ما مرقنا بسرعة جنونية جنب كم محطّة، قطعنا كم
\VUgras{نفق} \textsuperscript{(87)} وتركنا وراءنا \VUgras{بيوت}
\textsuperscript{(88)} \VUgras{حرّاس المزلقان} الزغيرة اللي ما إلها عدّ.
وهزّة القطار هدهدتني ونمت نوم عميق.

%%K t12-tit-4 sub
\VUsaut{3.0mm}
\VUpk{10.2pt}{40}{محطّة دويتش-آفريكور.}
\VUsaut{3.0mm}

%%K t12-19-1 p
19. --- وفجأة صحّيت على صوت موظّف عم يصرخ : « دويتش-آفريكور !
\textsuperscript{(*)} الكل ينزل ! » وصار تفتيش الجمرك وكل معاملات
الحدود المملّة. واضطرّيت ظبّط ساعة جيبتي على \VUgras{الساعة الكبيرة}
\textsuperscript{(89)} تبع المحطّة، اللي كانت مقدّمة خمسة وخمسين دقيقة ؛
وأنا بالكاد صاحي، صعب عليّ ميّز \VUgras{العقرب الكبير}
\textsuperscript{(90)} عن \VUgras{الزغير} \textsuperscript{(91)} ؛
و\VUgras{مينا الساعة} \textsuperscript{(92)} نفسها كانت مشوّشة تماماً من
ضباب الصبح.

%%K t12-20-1 p
20. --- ووقت موظّفي الجمرك عم يفتّشوا، كنت وأنا عم اتثاءب عم فكّ
\VUgras{الإعلانات} \textsuperscript{(93)} اللي مزيّنة حيطان المحطّة.
وفطرت لأمضّي الوقت، وتطلّعت بخريطة \VUgras{الشبكة} \textsuperscript{(94)}
الألمانية لشوف قدّيش لسّا في مسافة بيني وبين سترازبورغ، ورجعت طلعت
عالعربة، وأنا متقوّي بالأمل إنّي رح وصل عن قريب لهدف سفرتي.

\VUsaut{6.0mm}
\VUfilet{22mm}
